%%%%%%%%%%%%%%%%%%%%%%%%%%%%%%%%%%%%%%%%%%%%%%%%%%%%%%%%%%
%% Template for lab reppot
%% Antonin KEnzi
%% inspired by
%% template from reds institute for syslog 
%% and small repport by Olivier Morel
%%%%%%%%%%%%%%%%%%%%%%%%%%%%%%%%%%%%%%%%%%%%%%%%%%%%%%%%%%
\documentclass[11pt]{article}
\usepackage[T1]{fontenc}
\usepackage[utf8x]{inputenc} 
\usepackage[french]{babel} % Règles typographiques françaises, césures
% Ne pas toucher, nécessaire à la mise en forme
\usepackage{graphicx} % Insertion images
\usepackage{float}  % allow to force the placement of picture with 'H'
\usepackage[top=3cm, bottom=3cm, left=3cm, right=3cm, showframe=false]{geometry}
\usepackage{amsmath}
%\usepackage{hyperref}  % used to produce hyperlink when referencing but those links are marked by big red square on pdf softwares (in theory not when printed)
\usepackage{pdfpages}   % allow to include multipages pdf in document
\usepackage{multirow}
\usepackage{multicol}
\usepackage{lastpage}
\usepackage{datetime}
% More defined colors
\usepackage{xcolor}
%=================Police======================
%Les police suivante fonctionnent mais n'ont pas obligatoirement une equivalence "math"
%\usepackage[sfdefault]{inter} %% Option 'sfdefault' sans serif et pas de math mode
%\usepackage[]{mathpazo}

%\usepackage[math]{kurier} 
%\usepackage[italic]{mathastext}
\usepackage[sfdefault,scaled=1]{FiraSans}
\usepackage{amsmath,amsthm}
\usepackage[cmintegrals]{newtxsf}
%=================Package additionnel======================
% Pythontex------------------------------------------------$
%Permet d'excuter des script python
\usepackage{pythontex}
%% Configuration pour pythontex
\setpythontexlistingenv{pylisting}
\begin{pycode}
def SI(var, unit):
    return '\\SI{' + str(var) + '}{' + unit + '}'
def NUM(var, unit):
    return '\\num{' + str(var) + '}{' + unit + '}'
\end{pycode}


% siunitx--------------------------------------------
%Permet l'affichage des valeurs avec les normes SI
\usepackage{siunitx}	% also load some abbreviations
\sisetup{
	locale = FR,
	round-mode = figures,
	round-precision = 4,				% 4 digits in total
	scientific-notation = engineering,	% exponnant by multiple of 3
	per-mode            = symbol,
	per-symbol          = /,
	exponent-to-prefix  = true,
	output-decimal-marker = {.},		% decimal marker will always be printed as a '.'
	output-complex-root = j,			% imaginary unit will always be 'j'
	exponent-product = \cdot,			% exponant product will be a dot
	detect-all							% garde la même fonte que le reste du texte
}

% Tikz--------------------------------------------
%Permet la création de schémas/illustration en code
\usepackage{tikz}
\usetikzlibrary{positioning,calc,patterns,angles,quotes}
\newcommand{\gnufix}{gnuplot def/.append style={prefix={}}}

% circuitikz--------------------------------------------
%Permet la création de schémas en code (à vérifier l'implentation)
\usepackage[siunitx]{circuitikz}
\usetikzlibrary{positioning,calc,patterns,angles,quotes}

% minted--------------------------------------------
%Permet l'affichage de code dans les documents
\usepackage{minted}
%================= Configuration ==========================
%%%%%%%  Constantes (à changer) %%%%%%%%

% add any number of students
\newcommand{\nomAuteur}{Prénom \bsc{Nom} \& prénom \bsc{Nom}}

\newcommand{\titre}{Titre du labo}

\newcommand{\nbrLabo}{X} % numéro de laboratoire

\newcommand{\prof}{P.rof}

\newcommand{\assistant}{A.ssist}

\newcommand{\classe}{Classe}

\newcommand{\departement}{département}

\newcommand{\cours}{cours}

\newdate{date}{01}{06}{1996}

%caption pour equation : 
\newcommand{\capequation}[1]{\begin{center} \small #1 \end{center}}
%commande python à modifer pour que ça fonctionne correctement
\newcommand{\pySI}[2]{\py{'\\SI{' + str(#1) + '}{#2}'}}
\newcommand{\pyNUM}[2]{\py{'\\NUM[round-precision = #2]{' + str(#1) + '}'}}
% Pied de page---------------------------------------------
% Entête et pieds de page (A ADAPTER)
\pagestyle{empty}
\usepackage{fancyhdr}
\setlength{\headheight}{27.06pt}
\pagestyle{fancy}
% Entête :
\renewcommand{\headrulewidth}{1pt}
\fancyhead[L]{Labo \nbrLabo: \titre} 
\fancyhead[R]{Rapport de laboratoire} 
% Pied de page
\fancyfoot{} % clear all footer fields
\renewcommand{\footrulewidth}{1pt}
\fancyfoot[R]{Page \thepage / \pageref{LastPage}} 
\fancyfoot[L]{HEIG-VD | \nomAuteur }
%chemin des images 
\graphicspath{{graphics/}} %Setting the graphicspath
%=================== Script  =============================
\begin{document}
%%%%%%%%%%%%%%%%%%%%%%%%%%%%%%%%%%%%%%%%%%%%%%%%%%%%%%%%
% Ce document ne doit pas être modifié
% Il faut modifier les informations depuis le fichier
% constantes.tex
%%%%%%%%%%%%%%%%%%%%%%%%%%%%%%%%%%%%%%%%%%%%%%%%%%%%%%%%
\begin{center}
%\includegraphics[height=2cm]{./logo/logo_heig_IAI.png}
%\includegraphics[height=2cm]{./logo/logo_heig_reds.png}
\includegraphics[height=2cm]{./logo/logo_heig.pdf}
\hfill \includegraphics[height=2cm]{./logo/logo_hesso.pdf}
\end{center}

\vfill

\begin{center}
	\Large Laboratoire \nbrLabo:\\
	\huge \titre \\
\end{center}


\begin{center}
	\begin{tabular}{l l}
		Département :          & \departement\\
		Unité d'enseignement : & \cours\\
	\end{tabular}
\end{center}


\vfill

%\begin{center}
%	\includegraphics{Images/it.png}
%\end{center}

\begin{tabular}{l l}
	\Large	Auteurs :		& \Large \nomAuteur\\
	Professeur :			& \prof\\
	Assistant :				& \assistant\\
	Classe:					& \classe\\
	Date :					& \displaydate{date}\\
\end{tabular}

%\begin{center}
%\includegraphics[height=1.5cm]{./images/logo_reds.pdf}
%\end{center}
\thispagestyle{empty}
\setcounter{page}{1}	
% enlève le numéro sur la premiere page
\tableofcontents
\begin{thebibliography}{9}

\bibitem{label}
  name of,
  \textit{Tittle}, 
  where, when.

\end{thebibliography}

\end{document}